% Rencana Pembelajaran Semester (RPS)
% Berkas ini adalah MASUKAN contoh untuk pipeline AI Agent Buku Ajar.
%
% Sengaja ditulis sebagai teks biasa, bukan LaTeX yang harus dikompilasi:
% pada MVP ini RPS diperlakukan sebagai teks dan dikirim apa adanya ke Book
% Planner. Parser terstruktur CPMK/Sub-CPMK menyusul di Tahap 4 tanpa
% mengubah tipe BookRequest.

\section*{Identitas Mata Kuliah}

\begin{tabular}{ll}
Mata Kuliah   & Algoritma dan Struktur Data \\
Kode          & IF1234 \\
SKS           & 3 SKS (2 teori, 1 praktikum) \\
Semester      & 3 \\
Program Studi & Informatika \\
Prasyarat     & Dasar Pemrograman \\
Dosen         & Tim Pengampu \\
\end{tabular}

\section*{Deskripsi Mata Kuliah}

Mata kuliah ini membahas cara menyusun data agar operasi terhadapnya menjadi
efisien, serta teknik-teknik algoritmik dasar untuk mengolahnya. Penekanan
diberikan pada penalaran tentang biaya komputasi: mengapa satu rancangan
struktur data lebih tepat daripada yang lain untuk persoalan tertentu, dan
bagaimana memperkirakan biayanya sebelum menulis kode. Materi disampaikan
dengan pendekatan analisis kompleksitas yang diikuti praktik implementasi,
sehingga mahasiswa tidak sekadar menghafal struktur data tetapi mampu memilih
dan membenarkan pilihannya.

\section*{Capaian Pembelajaran Lulusan (CPL)}

\begin{enumerate}
  \item Mampu menerapkan pemikiran logis, kritis, sistematis, dan inovatif
        dalam konteks pengembangan ilmu pengetahuan dan teknologi.
  \item Mampu merancang, mengimplementasikan, dan menganalisis solusi
        komputasi yang efisien untuk persoalan nyata.
  \item Mampu mengomunikasikan gagasan teknis secara tertulis maupun lisan.
\end{enumerate}

\section*{Capaian Pembelajaran Mata Kuliah (CPMK)}

\begin{enumerate}
  \item CPMK-1: Mahasiswa mampu menganalisis kompleksitas waktu dan ruang
        sebuah algoritma menggunakan notasi asimtotik.
  \item CPMK-2: Mahasiswa mampu mengimplementasikan struktur data linear
        (senarai, tumpukan, antrean) dan menjelaskan konsekuensi biaya setiap
        operasinya.
  \item CPMK-3: Mahasiswa mampu menerapkan struktur data pohon, termasuk
        pohon pencarian biner dan pohon seimbang, untuk persoalan pencarian
        dan pengurutan.
  \item CPMK-4: Mahasiswa mampu memilih dan menerapkan teknik algoritmik
        dasar (pengurutan, pencarian, rekursi, pemrograman dinamis) sesuai
        karakteristik persoalan.
  \item CPMK-5: Mahasiswa mampu menilai trade-off antara waktu, ruang, dan
        kerumitan implementasi dalam memilih solusi.
\end{enumerate}

\section*{Sub-Capaian Pembelajaran Mata Kuliah (Sub-CPMK)}

\begin{enumerate}
  \item Sub-CPMK-1.1: Menjelaskan definisi O, Omega, dan Theta serta
        menghitungnya pada potongan kode sederhana.
  \item Sub-CPMK-2.1: Mengimplementasikan senarai berantai tunggal dan ganda
        beserta operasi penyisipan dan penghapusan.
  \item Sub-CPMK-2.2: Menjelaskan perbedaan perilaku tumpukan (LIFO) dan
        antrean (FIFO) serta menerapkannya pada persoalan nyata.
  \item Sub-CPMK-3.1: Menelusuri operasi pencarian dan penyisipan pada pohon
        pencarian biner, termasuk kasus terburuknya.
  \item Sub-CPMK-3.2: Menjelaskan mengapa pohon pencarian biner dapat menjadi
        tidak seimbang dan bagaimana pohon seimbang mengatasinya.
  \item Sub-CPMK-4.1: Membandingkan algoritma pengurutan dasar dan
        menjelaskan kapan masing-masing tepat dipakai.
  \item Sub-CPMK-4.2: Menerapkan rekursi dan pemrograman dinamis pada
        persoalan yang memiliki sub-persoalan bertumpang tindih.
  \item Sub-CPMK-5.1: Memberikan rekomendasi rancangan disertai alasan
        kuantitatif terhadap pilihan struktur data pada studi kasus.
\end{enumerate}

\section*{Rencana Mingguan}

\begin{enumerate}
  \item Minggu 1 - Kontrak kuliah; pengantar struktur data dan mengapa
        efisiensi diukur, bukan dirasakan.
  \item Minggu 2 - Analisis kompleksitas: notasi asimtotik, analisis kasus
        terburuk dan rata-rata, aturan penjumlahan dan perkalian.
  \item Minggu 3 - Senarai berantai: representasi simpul, senarai tunggal,
        senarai ganda, senarai melingkar; biaya penyisipan dan penghapusan.
  \item Minggu 4 - Tumpukan: operasi, implementasi berbasis larik dan
        senarai, penerapan pada evaluasi ekspresi dan penelusuran balik.
  \item Minggu 5 - Antrean: antrean biasa, antrean melingkar, dan antrean
        berprioritas; penerapan pada penjadwalan.
  \item Minggu 6 - Pengantar pohon: istilah, representasi, penelusuran
        preorder, inorder, dan postorder.
  \item Minggu 7 - Pohon pencarian biner: pencarian, penyisipan, penghapusan,
        dan kasus terburuk yang membuatnya linear.
  \item Minggu 8 - Ujian Tengah Semester.
  \item Minggu 9 - Pohon seimbang: AVL, rotasi, dan mengapa penyeimbangan
        menjaga biaya operasi tetap logaritmik.
  \item Minggu 10 - Antrean berprioritas dan tumpukan biner (heap):
        penyisipan, penghapusan, dan pengurutan dengan heap.
  \item Minggu 11 - Algoritma pengurutan dasar dan lanjut: gelembung,
        seleksi, sisipan, gabung, cepat, dan perbandingan biayanya.
  \item Minggu 12 - Tabel hash: fungsi hash, penanganan tabrakan dengan
        perantaian dan pengalamatan terbuka, faktor beban.
  \item Minggu 13 - Rekursi: relasi rekurens, kasus dasar versus kasus
        rekursif, dan biaya pemanggilan bertingkat.
  \item Minggu 14 - Pemrograman dinamis: memoization, penyelesaian
        bertahap, dan contoh persoalan klasik.
  \item Minggu 15 - Studi kasus terpadu: memilih struktur data untuk
        persoalan nyata dan membenarkan pilihannya secara kuantitatif.
  \item Minggu 16 - Ujian Akhir Semester.
\end{enumerate}

\section*{Metode Penilaian}

\begin{tabular}{lll}
Aktivitas              & Bobot & CPMK \\
Tugas dan kuis         & 20\%  & CPMK-1, CPMK-2 \\
Praktikum              & 25\%  & CPMK-2, CPMK-3, CPMK-4 \\
Ujian Tengah Semester  & 25\%  & CPMK-1, CPMK-2 \\
Ujian Akhir Semester   & 25\%  & CPMK-3, CPMK-4, CPMK-5 \\
Ujian Praktik          & 5\%   & CPMK-4 \\
\end{tabular}

\section*{Referensi}

\begin{enumerate}
  \item Cormen, T. H., Leiserson, C. E., Rivest, R. L., \& Stein, C.
        \emph{Introduction to Algorithms}. MIT Press.
  \item Sedgewick, R., \& Wayne, K. \emph{Algorithms}. Addison-Wesley.
  \item Weiss, M. A. \emph{Data Structures and Algorithm Analysis}.
        Pearson.
  \item Skiena, S. S. \emph{The Algorithm Design Manual}. Springer.
\end{enumerate}

\section*{Catatan Penyusunan Buku}

Buku ajar yang disusun dari RPS ini ditujukan untuk mahasiswa semester tiga
yang sudah menguasai dasar pemrograman prosedural. Setiap bab diharapkan
memuat tujuan pembelajaran yang eksplisit, minimal tiga contoh bertingkat
(dari yang langsung sampai yang menuntut penalaran), dan latihan yang
menutup materi. Bab tidak boleh mengandalkan rujukan yang tidak tersedia
secara terbuka, dan setiap angka kompleksitas yang disebut harus disertai
alasan mengapa angka itu yang berlaku.
